\chapter{Limitations and Threats to Validity}
\label{sec:limitations}

\paragraph{The scope of the comparison} is limited by two choices about the models.
First, both fine-tuned checkpoints are \gls{bart}-based, so the fine-tuned approach is represented by a single architecture family.
The obvious alternative is \gls{tfive} \parencite[1]{raffel2023exploringlimitstransferlearning}, which casts every text problem into one text-to-text format.
It was not trained or evaluated here, so the results say nothing about how the two architectures would compare.

Second, the prompted models were evaluated only in 4-bit quantised form, so that they could run locally on a laptop alongside the \gls{bart} checkpoints, as the browser extension's local deployment requires (\cref{sec:method-three}).
The effect of quantisation on their outputs was not measured.

\paragraph{Evaluation validity} is bounded by the benchmarks themselves.
WikiLarge's test split is a strict subset of ASSET's, so an improvement that holds on both confirms the scoring, not that the model generalises.
\qty{19.02}{\percent} of WikiLarge's validation split also appears in the training split.
Every validation loss quoted here is therefore a training diagnostic, not an estimate of generalisation.
Checkpoint selection is unaffected, because all checkpoints are validated on the same split.

ASSET's ten references were written under instructions for non-native speakers of English \parencite[13294]{kew-etal-2023-bless}, so metrics computed against them reward that style of simplification.
This matters most for the four other audiences, for which no reference data exist.

For roughly \qty{2}{\percent} of the document test split, no model could produce the target: the sources are disambiguation stubs, each paired with a full article as its reference.
Every system is scored against these references alike.

\paragraph{Metric validity} is the tighter constraint.
\Gls{lens} was trained on sentence-level human ratings \parencite[16385--16386]{maddela-etal-2023-lens} and is applied here to documents, out of its domain.
Its verdict therefore cannot fully settle its disagreement with \gls{dsari}.
It was not computed at sentence level, so the sentence-level comparison rests on \gls{ngram} and readability metrics alone.

\pagebreak

The \texttt{bert-score} implementation used here silently truncates inputs beyond roughly \num{510}~tokens, a second truncation point that need not coincide with the model's.
BERTScore also separates the fine-tuned and untuned document checkpoints by only \num{0.73}~points at \(\gls{nsize} = 500\) and \num{1.81} at \(\gls{nsize} = 8{,}000\), so it confirms nothing on its own.
Neither \gls{lens} nor BERTScore was significance-tested.
\Gls{dsari} here compares a corpus-trained system with a prompted one against a single reference, which is one reason the document-level ordering remains open.

\paragraph{Human judgement} was not collected.
All of these metrics are proxies for it, and \gls{bleu} and \gls{sari} correlate with human judgement only on some dimensions \parencite[148, 150--151]{alva-manchego-etal-2020-survey}.
They are complemented by structured manual inspection from the author.
No rating protocol was followed and no agreement statistics were computed.

Neither the metrics nor the inspection measure cohesion across sentences, so it remains open how far simplification disturbs it.
\Cref{sec:future-work} specifies both a cohesion measurement and a human study.

\paragraph{Confounds affect the headline comparisons} in four places.
First, the full-scope document run increased the corpus \num{6.6}-fold and the epoch count \num{2.5}-fold at once, so its gain cannot be attributed to either change alone.
The normaliser correction, by contrast, was measured separately and is not part of this confound (\cref{sec:document-results}).

Second, the two document runs differ in sample size.
The stability of the shared untuned checkpoint suggests that the comparison is informative, but it does not replace scoring both checkpoints on the same sample.
Third, scope and checkpoint also change together in the qualitative comparison.

Fourth, the comparison simplifier was fine-tuned with the prefix \texttt{Simplify:} on every source sentence, a preprocessing detail given in the training setup of the paper that released it \parencite[4]{cohen2026simplify}.
This project's evaluation and backend pass inputs without it.
With the prefix, its authors report \num{37.96}~\gls{sari} and \num{7.87}~\gls{fkgl} on the same \num{359} ASSET test sentences \parencite[3, 6]{cohen2026simplify}, against \num{38.17} and \num{7.85} measured here without it.
Both evaluations compute \gls{sari} and \gls{fkgl} with EASSE \parencite[12]{cohen2026simplify}.
The decoding differs: the profile they report for validation sets no repetition penalty \parencite[12]{cohen2026simplify}, whereas the one used here sets it to \num{1.2} (\cref{sec:finetuning}).
Their EASSE version is not reported either, so the two figures are not a controlled comparison.
They nonetheless give no sign that the omission lowered the comparison simplifier's score.

\pagebreak

\paragraph{Reproducibility} is bounded, and the bound was measured.
The numbers in this thesis were produced on four machines.
The same \num{500} documents, scored under CUDA and under Apple Silicon from two copies of the weights, produce byte-identical generations.
A rescore of cached generations on the \gls{cpu} reproduces a CUDA run bit-wise.
Deterministic beam-search output and the \gls{ngram} metrics computed from it.
Sampled output, which the prompted rows rely on, is not covered by this check.

\Gls{lens} runs only in a separate environment with a pinned library version.
The \gls{lens} column was therefore produced by a different software stack than the other metrics in the same table, and it cannot be recomputed in the service's environment.

The two unseeded training runs behind the noise floor (\cref{sec:sentence-results}) predate the provenance tooling and have no artefact.
They are quoted only as a variance observation.
There are two mojibake filters implementations, and every row count is quoted with the one it used.
Until automatic checkpoint pruning was switched off, training runs kept only their most recent checkpoints, so no per-epoch metric trajectory can be reconstructed for those early runs.

\paragraph{The guards and the content filter} rest on little labelled data.
The guards were not checked against labelled data, so no false-positive rate exists for any of them.
The measurements count how often they fire per request, since the response format does not allow a per-sentence rate (\cref{sec:deployment-results}).

The corpus deny-list rejects two strings, ``other websites'' and ``other pages'', but supports no claim about how often the checkpoints produce corpus boilerplate.
Prompt injection is not mitigated: the guards check only the generated output, and preventing it would require additional checks on the input.
The content filter's decisions, finally, were labelled for a sample of 69 items (\cref{sec:deployment-results}).

\paragraph{Deployment evaluation validity} is the weakest of these.
The site audit covers twelve pages, two per category, and reports no confidence intervals.
Both the site audit and the latency profiler run under \texttt{jsdom}, which by default does not run a page's own scripts \parencite{jsdom2026}.
The latency figures therefore do not include the time a browser needs to update and lay out the page.
Visual changes and pages rendered entirely by client-side JavaScript are not part of the audit; both were checked by hand on several pages during development.
The contention result comes from a single run.

The document-level error cases reported here come from the offline evaluation (\cref{sec:qualitative}).
The errors observed in the extension before the pipeline fixes reflect the pipeline rather than the model (\cref{tab:finding9}).
Finally, the document cut's coverage rate is counted but was not analysed, so it is unknown how much of each simplified section is left as the page wrote it.
